\chapter{Kapcsolodo projektanyagok}
\label{app:project-assets}

Ez a fuggelek a dolgozat keszitesehez felhasznalhato, mar meglevo projektanyagokat foglalja ossze. A dokumentumok nem a nyomtatott dolgozat reszei, de hasznos hatteret adnak a kovetelmenyek, a rendszerterv es a validacio kidolgozasahoz.

\section{Use case leirasok}

A \texttt{docs/usecases} konyvtar a legfontosabb felhasznaloi folyamatok leirasat tartalmazza. Ezek kulonosen hasznosak a kovetelmenyfejezet kialakitasahoz. A jelenlegi use case keszlet lefedi tobbek kozott a regisztraciot, a receptletrehozast, a keresest, a reszletes receptnezetet, a profilkezelest, valamint a kommenteleshez kapcsolodo folyamatokat.

\section{Wireframe-ek}

A \texttt{docs/wireframes} konyvtarban a fo oldalakhoz kapcsolodo SVG wireframe-ek talalhatok. Ezek jo kiindulasi pontot adnak a felhasznaloi felulet fejlodesenek bemutatasahoz. Ha a dolgozatba kepet szeretnel beemelni, erdemes a vegleges valtozatban PDF vagy PNG formatumban exportalt abrakat hasznalni.

\section{API leiras es tesztanyagok}

Az OpenAPI leiras a \texttt{docs/api/recipe\_api.yaml} allomanyban talalhato, mig a kezi kiprobalast segito Bruno gyujtemeny a \texttt{docs/bruno/api-calls} konyvtarban helyezkedik el. Ezek a fajlok a validacios fejezethez es az API szerzodes bemutatasahoz is hasznos alapanyagot szolgaltatnak.

